\makeatletter
\def\input@path{{../../core/}}
\makeatother

% Pengaturan ukuran teks
\documentclass[12pt,twoside]{report}

% Judul dokumen
\title{Buku Tugas Akhir ITS}
\author{Musk, Elon Reeve}

% Pengaturan ukuran halaman dan margin
% Untuk twoside: inner (binding side) = 30mm, outer (free side) = 20mm
\usepackage[a4paper,top=30mm,inner=30mm,outer=20mm,bottom=25mm,twoside]{geometry}

% Pengaturan ukuran spasi
\usepackage[singlespacing]{setspace}

% Pengaturan detail pada file PDF
\usepackage[bookmarksnumbered,pdfborder={0 0 0},hidelinks]{hyperref}
\hypersetup{pdftitle={Buku Tugas Akhir ITS}, pdfauthor={Musk, Elon Reeve}}

% Pengaturan pewarnaan
\usepackage[table,xcdraw]{xcolor}

% Pengaturan kutipan artikel
\usepackage[style=apa, backend=biber]{biblatex}
\PassOptionsToPackage{hyphens}{url}

% Package lainnya
\usepackage{changepage}
\usepackage{enumitem}
\usepackage{eso-pic}
\usepackage{fontspec}
\usepackage{etoolbox}
\usepackage{graphicx}
\usepackage{lipsum}
\usepackage{longtable}
\usepackage{tabularx}
\usepackage{wrapfig}
\usepackage{float}
\usepackage{algorithmic}
\usepackage{amsmath} % For enhanced math mode including \text command
\usepackage{amssymb} % For additional math symbols
\usepackage{titletoc} % Format daftar isi

% Sembunyikan subsection dari daftar isi
\setcounter{tocdepth}{1}

% Setup format daftar isi
\titlecontents{chapter}[0pt]
  {\fontsize{11}{13}\selectfont\rmfamily\mdseries}
  {BAB \thecontentslabel\hspace{1em}}
  {}
  {\hfill\contentspage}
  [\vspace{0pt}]

\titlecontents{section}[1.5em]
  {\fontsize{12}{14.4}\selectfont\rmfamily\mdseries}
  {\thecontentslabel\hspace{1em}}
  {}
  {\hfill\contentspage}
  [\vspace{0pt}]

\titlecontents{figure}[0pt]
  {\fontsize{11}{13}\selectfont\rmfamily\mdseries}
  {Gambar \thecontentslabel\hspace{1em}}
  {}
  {\hfill\contentspage}
  [\vspace{0pt}]

\titlecontents{table}[0pt]
  {\fontsize{11}{13}\selectfont\rmfamily\mdseries}
  {Tabel \thecontentslabel\hspace{1em}}
  {}
  {\hfill\contentspage}
  [\vspace{0pt}]

% Definisi untuk "Hati ini sengaja dikosongkan"
\patchcmd{\cleardoublepage}{\hbox{}}{
  \thispagestyle{empty}
  \vspace*{\fill}
  \begin{center}\textit{[Halaman ini sengaja dikosongkan]}\end{center}
  \vfill}{}{}

% Pengaturan penomoran halaman
\usepackage{fancyhdr}
\fancyhf{}
\renewcommand{\headrulewidth}{0pt}
\pagestyle{fancy}
\fancyfoot[LE,RO]{\thepage}
\patchcmd{\chapter}{plain}{fancy}{}{}
\patchcmd{\chapter}{empty}{plain}{}{}

% Command untuk bulan
\newcommand{\MONTH}{%
  \ifcase\the\month
  \or Januari% 1
  \or Februari% 2
  \or Maret% 3
  \or April% 4
  \or Mei% 5
  \or Juni% 6
  \or Juli% 7
  \or Agustus% 8
  \or September% 9
  \or Oktober% 10
  \or November% 11
  \or Desember% 12
  \fi}
\newcommand{\ENGMONTH}{%
  \ifcase\the\month
  \or January% 1
  \or February% 2
  \or March% 3
  \or April% 4
  \or May% 5
  \or June% 6
  \or July% 7
  \or August% 8
  \or September% 9
  \or October% 10
  \or November% 11
  \or December% 12
  \fi}

% Pengaturan format judul chapter
\usepackage{titlesec}
% Chapter (BAB): 14pt, bold, center | spacing before=0.21cm, after=0.64cm
\titleformat{\chapter}[hang]{\bfseries\fontsize{14}{16.8}\selectfont\centering}{BAB \arabic{chapter}}{1ex}{\centering}
% Section (Judul2): 13pt, bold, underline, left
\titleformat{\section}{\bfseries\fontsize{13}{15.6}\selectfont}{\thesection}{1em}{}
% Subsection (Judul3): 12pt, bold | spacing before=0.07cm, after=0cm
\titleformat{\subsection}{\bfseries\fontsize{12}{14.4}\selectfont}{\MakeUppercase{\thesubsection}}{1ex}{}
% Subsubsection (Judul4): 12pt, bold | spacing before=0.07cm, after=0cm
\titleformat{\subsubsection}{\bfseries\fontsize{12}{14.4}\selectfont}{\MakeUppercase{\thesubsubsection}}{1ex}{}
\titlespacing*{\chapter}{0pt}{0pt}{18pt}
\titlespacing{\section}{0pt}{12pt}{3pt}
\titlespacing{\subsection}{0pt}{6pt}{2pt}
\titlespacing{\subsubsection}{0pt}{6pt}{2pt}

% Pengaturan kedalaman penomoran section (x.x.x.x)
\setcounter{secnumdepth}{3}

% Atur variabel berikut sesuai namanya

% nama
\newcommand{\name}{Elon Reeve Musk}
\newcommand{\authorname}{Musk, Elon Reeve}
\newcommand{\nickname}{Elon}
\newcommand{\advisor}{Nikola Tesla, S.T., M.T}
\newcommand{\coadvisor}{Wernher von Braun, S.T., M.T}
\newcommand{\examinerone}{Dr. Galileo Galilei, S.T., M.Sc}
\newcommand{\examinertwo}{Friedrich Nietzsche, S.T., M.Sc}
\newcommand{\examinerthree}{Alan Turing, ST., MT}
\newcommand{\headofdepartment}{Prof. Albus Percival Wulfric Brian Dumbledore, S.T., M.T}

% identitas
\newcommand{\nrp}{0123 20 4000 0001}
\newcommand{\advisornip}{18560710 194301 1 001}
\newcommand{\coadvisornip}{18560710 194301 1 001}
\newcommand{\examineronenip}{18560710 194301 1 001}
\newcommand{\examinertwonip}{18560710 194301 1 001}
\newcommand{\examinerthreenip}{18560710 194301 1 001}
\newcommand{\headofdepartmentnip}{18810313 196901 1 001}

% judul
\newcommand{\tatitle}{Manajemen Sesi Agen LLM yang Stateful Berbasis Workspace dengan Menggunakan Snapshot MicroVM Firecracker}
\newcommand{\engtatitle}{Workspace-Centric Management of Stateful LLM Agent Sessions Using Snapshot-Based Firecracker MicroVM}

% kata kunci
\newcommand{\keywords}{microVM, Firecracker, Large Language Model, Model Context Protocol}

% tempat
\newcommand{\place}{Surabaya}

% jurusan
\newcommand{\studyprogram}{Teknologi Informasi}
\newcommand{\engstudyprogram}{Information Technology}

% fakultas
\newcommand{\faculty}{Teknik Elektro dan Informatika Cerdas}
\newcommand{\engfaculty}{Intelligent Electrical and Informatics Technology}

% singkatan fakultas
\newcommand{\facultyshort}{FTEIC}
\newcommand{\engfacultyshort}{ELECTICS}

% departemen
\newcommand{\department}{Teknologi Informasi}
\newcommand{\engdepartment}{Information Technology}

% kode mata kuliah
\newcommand{\coursecode}{TD1234}


\input{pustaka/tanda-hubung.tex}

% Menambahkan resource daftar pustaka
\addbibresource{../../core/pustaka/pustaka.bib}

% Pengaturan format potongan kode
\usepackage{listings}
\definecolor{comment}{RGB}{0,128,0}
\definecolor{string}{RGB}{255,0,0}
\definecolor{keyword}{RGB}{0,0,255}
\lstdefinestyle{codestyle}{
  commentstyle=\color{comment},
  stringstyle=\color{string},
  keywordstyle=\color{keyword},
  basicstyle=\footnotesize\ttfamily,
  numbers=left,
  numberstyle=\tiny,
  numbersep=5pt,
  xleftmargin=1em,
  frame=lines,
  breaklines=true,
  prebreak=\raisebox{0ex}[0ex][0ex]{\ensuremath{\hookleftarrow}},
  showstringspaces=false,
  upquote=true,
  tabsize=2,
  literate={"}{{\textquotedbl}}1,
}
\lstset{style=codestyle}

% Isi keseluruhan dokumen
\begin{document}

% Sampul Luar (Bahasa Indonesia)
\newcommand\covercontents{sampul/konten-id.tex}
\newcommand\coverleftmargin{30mm}
\AddToShipoutPictureBG*{
  \AtPageLowerLeft{
    % Ubah nilai berikut jika posisi horizontal background tidak sesuai
    \hspace{-4mm}

    % Ubah nilai berikut jika posisi vertikal background tidak sesuai
    \raisebox{0mm}{
      
\includegraphics[width=\paperwidth,height=\paperheight]{sampul/gambar/sampul-luar.png}
    }
  }
}

% Menyembunyikan nomor halaman
\thispagestyle{empty}

% Pengaturan margin untuk menyesuaikan konten sampul
\newgeometry{
  top=55mm,
  left=\coverleftmargin,
  right=20mm,
  bottom=20mm
}

\begin{flushleft}

  % Pemilihan font sans serif
  \sffamily

  % Pemilihan warna font putih
  \color{white}

  % Pemilihan font bold
  \fontseries{bx}
  \selectfont
  \begin{spacing}{1.5}
    \input{\covercontents}
  \end{spacing}

\end{flushleft}

\restoregeometry


% Sampul Luar (Bahasa Inggris)
\renewcommand\covercontents{sampul/konten-en.tex}
\renewcommand\coverleftmargin{20mm}
\AddToShipoutPictureBG*{
  \AtPageLowerLeft{
    % Ubah nilai berikut jika posisi horizontal background tidak sesuai
    \hspace{-4mm}

    % Ubah nilai berikut jika posisi vertikal background tidak sesuai
    \raisebox{0mm}{
      
\includegraphics[width=\paperwidth,height=\paperheight]{sampul/gambar/sampul-luar.png}
    }
  }
}

% Menyembunyikan nomor halaman
\thispagestyle{empty}

% Pengaturan margin untuk menyesuaikan konten sampul
\newgeometry{
  top=55mm,
  left=\coverleftmargin,
  right=20mm,
  bottom=20mm
}

\begin{flushleft}

  % Pemilihan font sans serif
  \sffamily

  % Pemilihan warna font putih
  \color{white}

  % Pemilihan font bold
  \fontseries{bx}
  \selectfont
  \begin{spacing}{1.5}
    \input{\covercontents}
  \end{spacing}

\end{flushleft}

\restoregeometry


% Sampul Luar Tipis (Bahasa Indonesia)
\renewcommand\covercontents{sampul/konten-id.tex}
\renewcommand\coverleftmargin{30mm}
\input{sampul/sampul-luar-tipis.tex}

% Sampul Luar Tipis (Bahasa Inggris)
\renewcommand\covercontents{sampul/konten-en.tex}
\renewcommand\coverleftmargin{20mm}
\input{sampul/sampul-luar-tipis.tex}

% Label tabel dan gambar dalam bahasa indonesia
\renewcommand{\figurename}{Gambar}
\renewcommand{\tablename}{Tabel}

% Pengaturan ukuran indentasi paragraf
\setlength{\parindent}{0.7cm}

% Pengaturan ukuran spasi paragraf
\setlength{\parskip}{1ex}

% Lembar pengesahan
\addcontentsline{toc}{chapter}{LEMBAR PENGESAHAN}
% Atur numbering halaman ke Roman dan mulai dari i
\pagenumbering{roman}
\setcounter{page}{1}

% Atur spacing untuk konten lembar pengesahan
\setlength{\parskip}{8pt}

% Macro untuk baris examiner/advisor (mengurangi redundansi)
\newcommand{\examrow}[3]{#1 & #2 \\ NIP: #3 & \\ & ................................... \\ & \\ & \\}

% Macro untuk baris dengan spacing (menghilangkan vspace repetitif)
\newcommand{\spline}[1]{\vspace{6pt} \\ #1}

\begin{center}
  \fontsize{14}{16.8}\selectfont
  \textbf{LEMBAR PENGESAHAN}
\end{center}
\vspace{10pt}

\begin{center}
  \textbf{\tatitle{}}
\end{center}

\begingroup
% Pemilihan font ukuran small
\small

\begin{center}
  \vspace{6pt}
  \textbf{TUGAS AKHIR} \\
  \vspace{6pt}
  Diajukan untuk memenuhi salah satu syarat \spline{memperoleh gelar Sarjana Komputer pada} \spline{Program Studi S-1 \studyprogram{}} \spline{Departemen \department{}} \spline{Fakultas \faculty{}} \spline{Institut Teknologi Sepuluh Nopember}
\end{center}

\begin{center}
  Oleh: \textbf{\name{}} \\[6pt]
   NRP. \nrp{}
\end{center}

\begin{center}
  Disetujui oleh Tim Penguji Tugas Akhir: 
\end{center}

\begingroup
% Menghilangkan padding
\setlength{\tabcolsep}{0pt}

\noindent
\begin{tabularx}{\textwidth}{X l}
  \examrow{\advisor{}}{\text{(Pembimbing I)}}{\advisornip{}}
  \examrow{\coadvisor{}}{\text{(Pembimbing II)}}{\coadvisornip{}}
  \examrow{\examinerone{}.}{\text{(Penguji I)}}{\examineronenip{}}
  \examrow{\examinertwo{}.}{\text{(Penguji II)}}{\examinertwonip{}}
  % \examrow{\examinerthree{}.}{\text{(Penguji III)}}{\examinerthreenip{}}
\end{tabularx}
\endgroup

\begin{center}
  \textbf{\MakeUppercase{\place{}}} \\
  \textbf{\MONTH{}}, \textbf{\the\year{}}
\end{center}
\endgroup


% \cleardoublepage

% Atur spacing untuk konten lembar pengesahan
\setlength{\parskip}{8pt}

\vspace{6pt}
\begin{center}
  \fontsize{14}{16.8}\selectfont
  \textbf{APPROVAL SHEET}
\end{center}
\vspace{18pt}

\begin{center}
  \textbf{\engtatitle{}}
\end{center}

\begingroup
% Pemilihan font ukuran small
\small

\begin{center}
  \vspace{6pt}
  \textbf{FINAL PROJECT} \\
  \vspace{6pt}
  Submitted to full one of the requirements \spline{for obtaining Bachelor degree at} \spline{Undergraduate Study Program of \engstudyprogram{}} \spline{Department of \engdepartment{}} \spline{Faculty of \engfaculty{}} \spline{Institut Teknologi Sepuluh Nopember}
\end{center}

\begin{center}
  By: \textbf{\name{}} \\[6pt]
  NRP. \nrp{}  
\end{center}

\begin{center}
  Approved by Final Project Examiner Team:
\end{center}

\begingroup
% Menghilangkan padding
\setlength{\tabcolsep}{0pt}

\noindent
\begin{tabularx}{\textwidth}{X l}
  \examrow{\advisor{}}{\text{(Advisor I)}}{\advisornip{}}
  \examrow{\coadvisor{}}{\text{(Co-Advisor II)}}{\coadvisornip{}}
  \examrow{\examinerone{}.}{\text{(Examiner I)}}{\examineronenip{}}
  \examrow{\examinertwo{}.}{\text{(Examiner II)}}{\examinertwonip{}}
  % \examrow{\examinerthree{}.}{\text{(Examiner III)}}{\examinerthreenip{}}
\end{tabularx}
\endgroup


\begin{center}
  \textbf{\MakeUppercase{\place{}}} \\
  \textbf{\ENGMONTH{}}, \textbf{\the\year{}}
\end{center}
\endgroup


\clearpage

% Halaman Pernyataan Orisinalitas
\addcontentsline{toc}{chapter}{PERNYATAAN ORISINALITAS}
% Atur numbering halaman ke Roman dan mulai dari i
\pagenumbering{roman}
\setcounter{page}{3}

\begin{center}
  \large
  \textbf{PERNYATAAN ORISINALITAS}
\end{center}

\vspace{2ex}

% Ubah paragraf-paragraf berikut sesuai dengan yang ingin diisi pada pernyataan keaslian

\noindent Yang bertanda tangan dibawah ini:

% hilangkan extra padding kiri
\noindent\begin{tabularx}{\textwidth}{@{}l l X@{}}
                       &   &                            \\
Nama Mahasiswa / NRP   & : & \nickname{} / \nrp{}           \\[6pt]
Departemen             & : & \department{}              \\[6pt]
Dosen Pembimbing 1 / NIP & : & \advisor{} / \advisornip{} \\[6pt]
Dosen Pembimbing 2 / NIP & : & \coadvisor{} / \coadvisornip{} \\
                       &   &                            \\
\end{tabularx}

\begin{sloppypar}
Dengan ini menyatakan bahwa Tugas Akhir dengan judul ``\tatitle{}'' adalah hasil karya sendiri, bersifat orisinal, dan ditulis dengan mengikuti kaidah penulisan ilmiah.
\end{sloppypar}

Bilamana di kemudian hari ditemukan ketidaksesuaian dengan pernyataan ini, maka saya bersedia menerima sanksi sesuai dengan ketentuan yang berlaku di Institut Teknologi Sepuluh Nopember.

\vspace{8ex}

\noindent\begin{tabularx}{\textwidth}{@{}>{\raggedright\arraybackslash}p{0.28\textwidth}X>{\raggedright\arraybackslash}p{0.28\textwidth}X>{\raggedright\arraybackslash}p{0.28\textwidth}@{}}
  & & & & \place{}, \MONTH{} \the\year{}               \\[6pt]
  Mengetahui,          & & & &            \\[6pt]
  Dosen Pembimbing 1   & & Dosen Pembimbing 2 & & Mahasiswa  \\
                       & & & &            \\
                       & & & &            \\
                       & & & &            \\
                       & & & &            \\
                       & & & &            \\
  \advisor{}           & & \coadvisor{}       & & \nickname{}    \\[20pt]
  NIP. \advisornip{}   & & NIP. \coadvisornip{} & & NRP. \nrp{} \\
\end{tabularx}

\clearpage

\begin{center}
  \large
  \textbf{STATEMENT OF ORIGINALITY}
\end{center}

\vspace{2ex}
% Ubah paragraf-paragraf berikut sesuai dengan yang ingin diisi pada pernyataan keaslian

\noindent The undersigned below:

\noindent\begin{tabularx}{\textwidth}{@{}l l X@{}}
                        &   &                            \\
  Name of student / NRP & : & \nickname{} / \nrp{}           \\[6pt]
  Department            & : & \engdepartment{}           \\[6pt]
  Advisor / NIP         & : & \advisor{} / \advisornip{} \\[6pt]
  Co-Advisor / NIP      & : & \coadvisor{} / \coadvisornip{} \\
                        &   &                            \\
\end{tabularx}

\begin{sloppypar}
Hereby declared that the Final Project with the title of ``\engtatitle{}'' is the result of my own work, is original, and is written by following the rules of scientific writing.
\end{sloppypar}

If in future there is a discrepancy with this statement, then I am willing to accept sanctions in accordance with provisions that apply at Sepuluh Nopember Institute of Technology.

\vspace{8ex}

\noindent\begin{tabularx}{\textwidth}{@{}>{\raggedright\arraybackslash}p{0.28\textwidth}X>{\raggedright\arraybackslash}p{0.28\textwidth}X>{\raggedright\arraybackslash}p{0.28\textwidth}@{}}
  & & & & \place{}, \ENGMONTH{} \the\year{}               \\[6pt]
  Acknowledged,          & & & &            \\[6pt]
  Advisor            & & Co-Advisor         & & Student    \\
                       & & & &            \\
                       & & & &            \\
                       & & & &            \\
                       & & & &            \\
                       & & & &            \\
  \advisor{}           & & \coadvisor{}       & & \nickname{}    \\[20pt]
  NIP. \advisornip{}   & & NIP. \coadvisornip{} & & NRP. \nrp{} \\
\end{tabularx}
\clearpage

% Nomor halaman pembuka dimulai dari sini
\pagenumbering{roman}

% Abstrak Bahasa Indonesia
% Atur numbering halaman ke Roman dan mulai dari i
\pagenumbering{roman}
\setcounter{page}{5}

\begin{center}
  {\fontsize{14}{16.8}\selectfont\textbf{ABSTRAK}}\\[6ex]
  {\fontsize{12}{14.4}\selectfont\textbf{\tatitle{}}} \\[4ex]
\end{center}

\addcontentsline{toc}{chapter}{ABSTRAK}

\vspace{2ex}

\begingroup
% Menghilangkan padding
\setlength{\tabcolsep}{0pt}

\noindent
\begin{tabularx}{\textwidth}{l >{\centering}m{2em} X}
  \textbf{Nama Mahasiswa / NRP}   & \textbf{:} & \textbf{\nickname{} / \nrp}        \\[6pt]
  \textbf{Departemen}             & \textbf{:} & \textbf{\department{} - Fakultas \faculty}      \\[20pt]
  \textbf{Dosen Pembimbing}       & \textbf{:} & \textbf{1. \advisor{}}   \\
                    &   & \textbf{2. \coadvisor{}} \\
\end{tabularx}
\endgroup

\noindent
\textbf{Abstrak}

% Ubah paragraf berikut dengan abstrak dari tugas akhir
Pada penelitian ini kami mengajukan \lipsum[1]

% Ubah kata-kata berikut dengan kata kunci dari tugas akhir
\vspace{2ex}
\noindent
\textbf{Kata Kunci: \keywords{}}

\clearpage

% Abstrak Bahasa Inggris
\begin{center}
  {\fontsize{14}{16.8}\selectfont\textbf{ABSTRACT}}\\[6ex]
  {\fontsize{12}{14.4}\selectfont\textbf{\engtatitle{}}} \\[4ex]
\end{center}

\addcontentsline{toc}{chapter}{ABSTRACT}

\vspace{2ex}

\begingroup
% Menghilangkan padding
\setlength{\tabcolsep}{0pt}

\noindent
\begin{tabularx}{\textwidth}{l >{\centering}m{2em} X}
  \textbf{Student Name / NRP}   & \textbf{:} & \textbf{\nickname{} / \nrp{}}        \\[6pt]
  \textbf{Department}           & \textbf{:} & \textbf{\engdepartment{} - Faculty of \engfaculty}      \\[20pt]
  \textbf{Advisor}              & \textbf{:} & \textbf{1. \advisor{}}   \\
                                &             & \textbf{2. \coadvisor{}} \\
\end{tabularx}
\endgroup

\noindent
\textbf{Abstract}

% Ubah paragraf berikut dengan abstrak dari tugas akhir
Pada penelitian ini kami mengajukan \lipsum[1]

% Ubah kata-kata berikut dengan kata kunci dari tugas akhir
\vspace{2ex}
\noindent
\textbf{Keywords: \keywords{}}

\clearpage

% Daftar isi
\addcontentsline{toc}{chapter}{DAFTAR ISI}
\begin{center}
  {\fontsize{14}{14.4}\selectfont\bfseries DAFTAR ISI}
\end{center}
\vspace{2ex}
\makeatletter
\@starttoc{toc}
\makeatother
\clearpage

% Daftar gambar
\addcontentsline{toc}{chapter}{DAFTAR GAMBAR}
\begin{center}
  {\fontsize{14}{14.4}\selectfont\bfseries DAFTAR GAMBAR}
\end{center}
\begingroup
\setlength{\parskip}{0pt}
\vspace{2ex}
\makeatletter
\@starttoc{lof}
\makeatother
\endgroup
\clearpage

% Daftar tabel
\addcontentsline{toc}{chapter}{DAFTAR TABEL}
\begin{center}
  {\fontsize{14}{14.4}\selectfont\bfseries DAFTAR TABEL}
\end{center}
\begingroup
\vspace{2ex}
\setlength{\parskip}{0pt}
\makeatletter
\@starttoc{lot}
\makeatother
\endgroup
\clearpage

% Nomor halaman isi dimulai dari sini
\pagenumbering{arabic}

% Margin dan spacing untuk konten bab
% Untuk twoside: inner (binding side) = 25mm, outer (free side) = 25mm
\newgeometry{top=20mm,inner=25mm,outer=25mm,bottom=25mm}
\setlength{\parskip}{8pt}

% Bab 1 pendahuluan
\input{chapter/1-pendahuluan.tex}
\cleardoublepage

% Bab 2 tinjauan pustaka
\chapter{TINJAUAN PUSTAKA}
\label{chap:tinjauan-pustaka}

% Ubah bagian-bagian berikut dengan isi dari tinjauan pustaka

Demi mendukung penelitian ini, \lipsum[1][1-5]

\section{Hasil Penelitian Terdahulu}

Beberapa penelitian terdahulu yang relevan dengan topik ini antara lain \lipsum[1][1-5]

\setlength{\LTpre}{6pt}
\setlength{\LTpost}{6pt}
\begin{longtable}{|c|p{4cm}|p{8cm}|}
  \caption{Penelitian Terdahulu}\label{tab:penelitian-terdahulu} \\
  \hline
  \textbf{No} & \textbf{Penelitian} & \textbf{Pembahasan} \\
  \hline
  \endfirsthead

  \multicolumn{3}{c}{\tablename\ \thetable\ -- \textit{Lanjutan dari halaman sebelumnya}} \\
  \hline
  \textbf{No} & \textbf{Penelitian} & \textbf{Pembahasan} \\
  \hline
  \endhead

  \hline
  \endfoot

  \hline
  \endlastfoot

  1 & Voyager & \parencite{wang2023voyager} \lipsum[2][1-3] \\
  \hline
  2 & Model Context Protocol & \parencite{cilacap2024mcp} \lipsum[2][1-3] \\
  \hline

\end{longtable}

\section{Roket Luar Angkasa}
\label{sec:roket-luar-angkasa}

% Contoh input gambar
\begin{figure}[H]
  \centering

  % Ubah dengan nama file gambar dan ukuran yang akan digunakan
  \includegraphics[scale=0.35]{gambar/roketluarangkasa.jpg}

  % Ubah dengan keterangan gambar yang diinginkan
  \caption{Peluncuran roket luar angkasa \emph{Discovery} \parencite{roketluarangkasa}.}
  \label{fig:roket-luar-angkasa}
\end{figure}

Roket luar angkasa merupakan \lipsum[1]

\emph{Discovery}, Gambar \ref{fig:roket-luar-angkasa}, merupakan \lipsum[2]

\section{Gravitasi}
\label{sec:gravitasi}

Gravitasi merupakan \lipsum[1]

\subsection{Hukum Newton}
\label{subsec:hukum-newton}

Newton \parencite{newton1687} pernah merumuskan bahwa \lipsum[1]
Kemudian menjadi persamaan seperti pada persamaan \ref{eq:hukum-pertama-newton}.

% Contoh pembuatan persamaan
\begin{equation}
  \label{eq:hukum-pertama-newton}
  \sum \mathbf{F} = 0\; \Leftrightarrow\; \frac{\mathrm{d} \mathbf{v} }{\mathrm{d}t} = 0.
\end{equation}

\subsection{Anti Gravitasi}
\label{subsec:anti-gravitasi}

Anti gravitasi merupakan \lipsum[1]

\cleardoublepage

% Bab 3 metodologi
\input{chapter/3-metodologi.tex}
\cleardoublepage

% Bab 4 hasil dan pembahasan
\input{chapter/4-hasil-pembahasan.tex}
\cleardoublepage

% Bab 5 kesimpulan dan saran
\input{chapter/5-kesimpulan-saran.tex}
\cleardoublepage

\chapter*{DAFTAR PUSTAKA}
\addcontentsline{toc}{chapter}{DAFTAR PUSTAKA}
\renewcommand\refname{}
\vspace{2ex}
\renewcommand{\bibname}{}
\begingroup
\def\chapter*#1{}
\printbibliography
\endgroup
\cleardoublepage

% Biografi penulis
\input{lainnya/biografi-penulis.tex}
\cleardoublepage

\end{document}
